\subsection{Duality of Z/nZ-modules}

In this No.~we denote by n an integer \textgreater0 and by T a cyclic group of order n.~A group G is said to be annihilated by n if \(g^{n}=1\) for all \(g\in G\) ; if moreover G is commutative, the group structure of G is underlying a unique Z/nZ-module structure.

For every group G we denote by \(\operatorname{Hom}(G, T)\) the group of homomorphisms of G into T; it is a commutative group annihilated by n.

PROPOSITION 7. --- Let G be a commutative group annihilated by n and H a subgroup of G. Then the restriction homomorphism \(\operatorname{Hom}(G, T) \to \operatorname{Hom}(H, T)\) is surjective.

For let \(f:H\to T\) be a homomorphism; we shall prove that there exists a homomorphism of G into T extending f.~Suppose first that G is cyclic, generated by an element g of order r dividing n; denote by t a generator of T. There exists a divisor s of r such that H is generated by \(g^{s}\) (I, p.~50, Prop. 19), and for every \(x \in Z\) we have \(f(g^{sx}) = t^{ax}\) , where a is an integer such that n divides ar/s. Then \(a/s = (ar/ns)(n/r)\) is an integer and the homomorphism \(g^{x} \mapsto t^{(a/s)x}\) , \(x \in Z\) , of G into T extends f.~In the general case consider the set of pairs \((\mathrm{H}', f')\) where \(H'\) is a subgroup of G containing H and \(f'\) is a homomorphism of \(H'\) into T extending f, and let us order this set by the relation \((\mathrm{H}', f') \leqslant (\mathrm{H}'', f'')\) if \(H' \subset H''\) and the restriction of \(f''\) to \(H'\) is \(f'\) . By Set Theory, III, p.~154, Def. 3 and Th. 2, this set has a maximal element \((\mathrm{H}_{1}, f_{1})\) and it suffices to prove that \(H_{1} = G\) ; if this is not the case, there exists \(g \in G\) , \(g \notin H_{1}\) and it is enough to prove that \(f_{1}\) may be extended to a homomorphism T of the subgroup of G generated by \(H_{1}\) and g. Now if C denotes the cyclic group generated by g, the restriction of \(f_{1}\) to \(C \cap H_{1}\) extends to a homomorphism \(f_{2}\) of C into T and the homomorphism \(xy \mapsto f_{1}(x) f_{2}(y)\) , \(x \in H_{1}\) , \(y \in C\) , of \(H_{1}C\) into T is the desired mapping.

COROLLARY 1. --- If G is a commutative group annihilated by n, and \(G \neq \{1\}\) , then \(\operatorname{Hom}(G, T) \neq \{1\}\) .

For it suffices to note that if H is a cyclic subgroup of G distinct from \(\{1\}\) , then \(\operatorname{Hom}(H, T) \neq \{1\}\) , and to apply Prop. 7.

COROLLARY 2. --- If G is a finite commutative group annihilated by n, then the groups G and Hom(G, T) have the same order.

If G is cyclic of order r, with generator g, then the mapping \(f \mapsto f(g)\) is a bijection of \(\operatorname{Hom}(G, T)\) onto the set of elements t of T such that \(t^{r} = 1\) , hence the assertion follows in this case. On the other hand, if H is a cyclic subgroup of G, we have \(\operatorname{Card}(G) = \operatorname{Card}(H)\) . \(\operatorname{Card}(G/H)\) ; further, we have an exact sequence

\[
\{1 \} \rightarrow \operatorname{Hom} (\mathrm{G} / \mathrm{H}, \mathrm{T}) \rightarrow \operatorname{Hom} (\mathrm{G}, \mathrm{T}) \rightarrow \operatorname{Hom} (\mathrm{H}, \mathrm{T}) \rightarrow \{1 \}
\]

(II, p.~227, Th. 1 and Prop. 7 above), hence

\[
\operatorname{Card} (\operatorname{Hom} (G, T)) = \operatorname{Card} (\operatorname{Hom} (H, T)). \operatorname{Card} (\operatorname{Hom} (G / H, T)).
\]

Since \(\text{Card}(\text{Hom}(H, T)) = \text{Card}(H)\) , it comes to the same to prove the corollary for G or for G/H. Now the result follows by induction on \(\text{Card}(G)\) .

Now let G and H be two groups and \(f: G \times H \to T\) a bimultiplicative mapping, that is, such that for any \(g, g' \in G\) , \(h, h' \in H\) we have

\[
f (g g ^ {\prime}, h) = f (g, h) f (g ^ {\prime}, h), \quad f (g, h h ^ {\prime}) = f (g, h) f (g, h ^ {\prime}).
\]

We define group homomorphisms

\[
s _ {f} \colon \mathrm{G} \to \operatorname{Hom} (\mathrm{H}, \mathrm{T}), d _ {f} \colon \mathrm{H} \to \operatorname{Hom} (\mathrm{G}, \mathrm{T}),
\]

by \(s_f(g)(h) = d_f(h)(g) = f(g, h)\) (cf.~II, p.~268, Cor. to Prop. 1, when G and H are commutative).

PROPOSITION 8. --- Suppose that G and H are commutative and annihilated by n.~If \(s_{f}\) is bijective and H is finite, then \(d_{f}\) is bijective and we have \(\text{Card}(G) = \text{Card}(H)\) .

If \(s_{f}\) is bijective and H is finite, we have by Cor. 2 to Prop. 7 the relation \(\operatorname{Card}(G) = \operatorname{Card}(\operatorname{Hom}(H, T)) = \operatorname{Card}(H)\) , hence \(\operatorname{Card}(G)\) is finite and by another application of the Corollary, \(\operatorname{Card}(\operatorname{Hom}(G, T)) = \operatorname{Card}(H)\) . So it is enough to prove that \(d_{f}\) is injective. Now if \(h \in \operatorname{Ker}(d_{f})\) , we have \(f(g, h) = 1\) for all \(g \in G\) , hence since \(s_{f}\) is bijective, \(\varphi(h) = 1\) for all \(\varphi \in \operatorname{Hom}(H, T)\) ; by Prop. 7 this implies \(\operatorname{Hom}(\operatorname{Ker}(d_{f}), T) = \{1\}\) , hence \(\operatorname{Ker}(d_{f}) = \{1\}\) by Cor. 1 to Prop. 7.

